\chapter{Geometria Proiettiva, Omografie e Modellizzazione della Telecamera}

\section{Richiami sullo Spazio Proiettivo e Basi Proiettive}
I modelli geometrici classici (come lo spazio affine) non consentono di modellare adeguatamente le proiezioni prospettiche tipiche dei sistemi di visione[cite: 1]. Per superare questo limite viene introdotto lo spazio proiettivo $\mathbb{P}^n$, costruito a partire dallo spazio vettoriale $\mathbb{R}^{n+1} \setminus \{\mathbf{0}\}$ definendo una relazione di equivalenza $\mathcal{R}$[cite: 1]:
\begin{equation}
\mathbf{x} \sim \mathbf{y} \iff \exists \lambda \neq 0 \quad \text{tale che} \quad \mathbf{x} = \lambda \mathbf{y}
\end{equation}
Lo spazio proiettivo $n$-dimensionale $\mathbb{P}^n$ è lo spazio quoziente delle classi di equivalenza rispetto a tale relazione[cite: 1].

\subsection{Base Proiettiva}
Una base per uno spazio proiettivo di dimensione $n$ ($\mathbb{P}^n$) è composta da $n+2$ punti proiettivi[cite: 1]:
\begin{itemize}
    \item $n+1$ punti i cui vettori rappresentanti sono linearmente indipendenti in $\mathbb{R}^{n+1}$[cite: 1].
    \item $1$ punto unità (o di calibrazione) necessario per fissare i fattori di scala dei rappresentanti[cite: 1].
\end{itemize}

Esempi di dimensioni di basi nei diversi spazi[cite: 1]:
\begin{itemize}
    \item \textbf{Retta proiettiva ($\mathbb{P}^1$):} La base richiede $3$ punti[cite: 1].
    \item \textbf{Piano proiettivo ($\mathbb{P}^2$):} La base richiede $4$ punti[cite: 1].
\end{itemize}

\subsection{Confronto tra Trasformazioni Affini e Proiettive}
Nelle trasformazioni affini sul piano, il parallelismo viene preservato[cite: 1]. Pertanto, conoscendo l'immagine di $3$ punti (es. i tre vertici di un quadrato), il quarto punto è vincolato ed è univocamente determinabile tramite l'intersezione delle rette parallele[cite: 1].

Nelle trasformazioni proiettive, il parallelismo \textbf{non} viene conservato[cite: 1]. Di conseguenza, l'immagine di $3$ punti non è sufficiente a definire la mappa, poiché il quarto punto può trovarsi in una posizione arbitraria[cite: 1]. Sono quindi indispensabili $4$ punti nel piano per identificare univocamente una trasformazione proiettiva (corrispondenti esattamete ai $4$ punti della base proiettiva di $\mathbb{P}^2$)[cite: 1].

\section{Immersione dello Spazio Affine e Punti all'Infinito}
L'immersione dello spazio affine $n$-dimensionale $\mathbb{R}^n$ nello spazio proiettivo $\mathbb{P}^n$ avviene esprimendo i punti in coordinate omogenee $\mathbf{x} = (x_1, x_2, \dots, x_n, x_{n+1})^T$[cite: 1]:
\begin{enumerate}
    \item \textbf{Punti al finito:} Se $x_{n+1} \neq 0$, è sempre possibile dividere tutte le coordinate per $x_{n+1}$, ottenendo un rappresentante normalizzato con l'ultima coordinata pari ad $1$: $(x_1/x_{n+1}, \dots, x_n/x_{n+1}, 1)^T$[cite: 1].
    \item \textbf{Punti all'infinito (o ideali):} Se $x_{n+1} = 0$, il punto ha forma $(x_1, x_2, \dots, x_n, 0)^T$[cite: 1]. Tali punti rappresentano le direzioni delle rette nello spazio e non appartengono allo spazio affine originale[cite: 1].
\end{enumerate}

\section{Dualità nel Piano Proiettivo $\mathbb{P}^2$}
In $\mathbb{R}^2$, l'equazione implicita di una retta è $ax + by + c = 0$[cite: 1]. Esprimendo il punto in coordinate omogenee $\mathbf{x} = (x_1, x_2, x_3)^T$, la retta equivale alla condizione di ortogonalità[cite: 1]:
\begin{equation}
\mathbf{l}^T \mathbf{x} = 0 \quad \text{con} \quad \mathbf{l} = \begin{pmatrix} a \\ b \\ c \end{pmatrix}
\end{equation}
Poiché i parametri $(a, b, c)$ identificano la stessa retta a meno di un fattore di scala non nullo ($\lambda a x + \lambda b y + \lambda c = 0$), il vettore $\mathbf{l}$ appartiene a sua volta a $\mathbb{P}^2$[cite: 1].

Esiste quindi una perfetta \textbf{dualità} nel piano proiettivo: sia i punti sia le rette sono rappresentati da vettori a $3$ componenti in $\mathbb{P}^2$, e le loro definizioni algebraiche sono interscambiabili[cite: 1]:
\begin{itemize}
    \item \textbf{Retta per due punti:} La retta $\mathbf{l}$ passante per i punti $\mathbf{p}$ e $\mathbf{q}$ è data dal prodotto vettoriale:
    \begin{equation}
    \mathbf{l} = \mathbf{p} \times \mathbf{q}
    \end{equation}
    \item \textbf{Intersezione di due rette:} Il punto di intersezione $\mathbf{x}$ tra due rette $\mathbf{l}_1$ e $\mathbf{l}_2$ è dato dal prodotto vettoriale:
    \begin{equation}
    \mathbf{x} = \mathbf{l}_1 \times \mathbf{l}_2
    \end{equation}
\end{itemize}

\subsection{La Retta all'Infinito}
In $\mathbb{P}^2$, la condizione $x_3 = 0$ individua l'insieme di tutti i punti all'infinito[cite: 1]. Tali punti giacciono su un'unica retta proiettiva, detta \textbf{retta all'infinito} $\mathbf{l}_\infty$[cite: 1]:
\begin{equation}
\mathbf{l}_\infty = \begin{pmatrix} 0 \\ 0 \\ 1 \end{pmatrix}
\end{equation}

\section{Gerarchia delle Trasformazioni Geometriche e Gradi di Libertà}
Le trasformazioni da $\mathbb{P}^n$ in $\mathbb{P}^n$ sono rappresentate da matrici quadrate ordinarie $(n+1) \times (n+1)$ definite a meno di un fattore di scala[cite: 1]. La struttura a blocchi generale di una matrice di trasformazione è:
\begin{equation}
H = \begin{bmatrix} A_{n \times n} & \mathbf{t}_{n \times 1} \\ \mathbf{v}^T_{1 \times n} & v \end{bmatrix}
\end{equation}

Di seguito viene descritta la gerarchia delle trasformazioni, ordinate dal livello più vincolato a quello più generale:

\begin{enumerate}
    \item \textbf{Isometrie (Euclidee):}
    \begin{itemize}
        \item Struttura: $A = R$ (matrice di rotazione ortogonale), $\mathbf{v} = \mathbf{0}$, $v = 1$[cite: 1].
        \item Invarianti: Distanze, angoli, aree[cite: 1].
        \item Gradi di Libertà (DoF):
        \begin{itemize}
            \item In 2D ($\mathbb{P}^2$): $3$ DoF ($1$ angolo di rotazione + $2$ componenti di traslazione)[cite: 1].
            \item In 3D ($\mathbb{P}^3$): $6$ DoF ($3$ angoli di Eulero + $3$ componenti di traslazione)[cite: 1].
        \end{itemize}
    \end{itemize}

    \item \textbf{Similarità:}
    \begin{itemize}
        \item Struttura: $A = sR$ (rotazione scalata di un fattore $s$), $\mathbf{v} = \mathbf{0}$, $v = 1$[cite: 1].
        \item Invarianti: Rapporti tra distanze, angoli, forme[cite: 1].
        \item Gradi di Libertà (DoF):
        \begin{itemize}
            \item In 2D ($\mathbb{P}^2$): $4$ DoF ($3$ delle isometrie + $1$ fattore di scala $s$)[cite: 1].
            \item In 3D ($\mathbb{P}^3$): $7$ DoF ($6$ delle isometrie + $1$ fattore di scala $s$)[cite: 1].
        \end{itemize}
    \end{itemize}

    \item \textbf{Trasformazioni Affini:}
    \begin{itemize}
        \item Struttura: $A$ matrice non singolare arbitraria, $\mathbf{v} = \mathbf{0}$, $v = 1$[cite: 1].
        \item Invarianti: Parallelismo, rapporti tra aree, retta all'infinito (mappa punti al finito in punti al finito)[cite: 1].
        \item Gradi di Libertà (DoF):
        \begin{itemize}
            \item In 2D ($\mathbb{P}^2$): $6$ DoF ($4$ elementi di $A_{2 \times 2}$ + $2$ di $\mathbf{t}$)[cite: 1].
            \item In 3D ($\mathbb{P}^3$): $12$ DoF ($9$ elementi di $A_{3 \times 3}$ + $3$ di $\mathbf{t}$)[cite: 1].
        \end{itemize}
    \end{itemize}

    \item \textbf{Trasformazioni Proiettive (Omografie):}
    \begin{itemize}
        \item Struttura: Matrice $Q_{(n+1) \times (n+1)}$ completamente generica (nessun vincolo su $\mathbf{v}$ e $v$)[cite: 1].
        \item Invarianti: Allineamento di punti, \emph{birapporto} (\emph{cross-ratio})[cite: 1]. Mappa punti al finito anche in punti all'infinito e viceversa[cite: 1].
        \item Gradi di Libertà (DoF): Poiché la matrice è definita a meno di un fattore di scala globale, si sottrae $1$ DoF dal numero totale di elementi[cite: 1]:
        \begin{itemize}
            \item In 2D ($\mathbb{P}^2 \to \mathbb{P}^2$): $3 \times 3 - 1 = 8$ DoF[cite: 1].
            \item In 3D ($\mathbb{P}^3 \to \mathbb{P}^3$): $4 \times 4 - 1 = 15$ DoF[cite: 1].
        \end{itemize}
    \end{itemize}
\end{enumerate}

\section{Omografie e Applicazioni Pratiche (Panorami)}
Un'omografia è un isomorfismo dello spazio proiettivo in sé ($\mathbb{P}^2 \to \mathbb{P}^2$), che rappresenta una trasformazione lineare, biunivoca e invertibile tra due piani proiettivi[cite: 1].

\subsection{Generazione di Viste Panoramiche}
Quando una telecamera ruota attorno al proprio centro ottico o inquadra una scena perfettamente piana/molto lontana, le variazioni di profondità della scena diventano trascurabili[cite: 1]. In questo caso:
\begin{itemize}
    \item La proiezione geometrica tra le diverse immagini scattate può essere modellata come un'omografia pianare invertibile[cite: 1].
    \item Tutte le immagini possono essere riproiettate ed allineate su un unico piano immagine comune, consentendo la creazione di mosaici e viste panoramiche[cite: 1].
\end{itemize}

Se il medesimo processo viene applicato ad oggetti vicini con variazioni di profondità significative, il modello piano fallisce: la trasformazione reale non è più invertibile e il mosaico risulterà fortemente distorto ed irrealistico[cite: 1].

\section{Modello della Telecamera Prospettica e Calibrazione}

\subsection{Proiezione Prospettica Ideale}
Sia $O$ il centro di proiezione (centro ottico) e $\Pi$ il piano immagine posto a distanza focale $f$ lungo l'asse ottico $Z$[cite: 1]. Un punto 3D $P = (X, Y, Z)^T$ viene proiettato sul piano immagine nel punto $p = (x, y)^T$ secondo le equazioni non lineari[cite: 1]:
\begin{equation}
x = f \frac{X}{Z}, \quad y = f \frac{Y}{Z}
\end{equation}
L'intersezione dell'asse ottico $Z$ con il piano immagine $\Pi$ definisce il \textbf{punto principale} $O_p = (x_0, y_0)$[cite: 1].

\subsection{Parametri Intrinseci e Coordinate Pixel}
Nelle telecamere reali, il sistema di riferimento dell'immagine ha l'origine nell'angolo in alto a sinistra del sensore e le coordinate sono misurate discrezionalmente in pixel[cite: 1]. Occorre considerare:
\begin{itemize}
    \item $k_x, k_y$: Numero di pixel per unità di lunghezza nelle direzioni orizzontale e verticale[cite: 1].
    \item $f_x = f \cdot k_x, \; f_y = f \cdot k_y$: Distanze focali espresse direttamente in unità pixel[cite: 1].
    \item $(x_0, y_0)$: Coordinate del punto principale espressa nel sistema di riferimento dei pixel[cite: 1].
    \item $\theta$: Angolo tra gli assi dei pixel (se gli elementi sensibili non sono perfettamente ortogonali; idealmente $\theta = 90^\circ$)[cite: 1].
\end{itemize}

La matrice dei \textbf{parametri intrinseci} $K$ ($3 \times 3$), assumendo $\theta = 90^\circ$, assume la forma[cite: 1]:
\begin{equation}
K = \begin{bmatrix} f_x & 0 & x_0 \\ 0 & f_y & y_0 \\ 0 & 0 & 1 \end{bmatrix}
\end{equation}
La matrice $K$ contiene $5$ parametri intrinseci liberi ($f_x, f_y, x_0, y_0$ ed eventualmente lo skew $s$)[cite: 1].

\subsection{Parametri Estrinseci e Matrice di Proiezione Completa}
I punti del mondo reale sono espressi in un sistema di riferimento arbitrario (Mondo)[cite: 1]. Per portarli nel sistema di riferimento della telecamera, si applica una rotazione $R$ ($3 \times 3$) e una traslazione $\mathbf{t}$ ($3 \times 1$), dette \textbf{parametri estrinseci} ($6$ DoF)[cite: 1].

Combinando parametri intrinseci ed estrinseci, la proiezione di un punto 3D espresso in coordinate omogenee $\mathbf{X}_W = (X, Y, Z, 1)^T$ nel punto pixel $\mathbf{x} = (u, v, 1)^T$ diventa una \textbf{trasformazione lineare proiettiva} regolata dalla matrice di proiezione $P$ ($3 \times 4$)[cite: 1]:
\begin{equation}
\mathbf{x} \sim P \mathbf{X}_W = K \begin{bmatrix} R & \mathbf{t} \end{bmatrix} \mathbf{X}_W
\end{equation}

\subsection{Conteggio complessivo dei Gradi di Libertà}
La matrice di proiezione $P$ ha dimensione $3 \times 4$ ($12$ elementi) ed è definita a meno di un fattore di scala globale, fornendo $11$ gradi di libertà totali[cite: 1]:
\begin{equation}
11 \text{ DoF totali} = 5 \text{ parametri intrinseci } (K) + 6 \text{ parametri estrinseci } (R, \mathbf{t})
\end{equation}
Il processo di \textbf{calibrazione della telecamera} consiste nello stimare numericamente i parametri della matrice $P$ (e di conseguenza disaccoppiare $K$, $R$ e $\mathbf{t}$) risolvendo opportuni sistemi lineari basati su punti di corrispondenza noti[cite: 1].